\section{MAXIMAL SUBALGEBRAS OF SEMI-SIMPLE LIE ALGEBRAS}

THEOREM 1. Let V be a finite dimensional vector space, g a reductive Lie subalgebra in \(\mathfrak{gl}(V)\) , q a Lie subalgebra of g and \(\Phi\) the bilinear form \((x,y)\mapsto\mathrm{Tr}(xy)\) on \(g\times g\) . Assume that the orthogonal complement n of q with respect to \(\Phi\) is a Lie subalgebra of g consisting of nilpotent endomorphisms of V. Then q is a parabolic subalgebra of g.

\begin{enumerate}
\def\labelenumi{\alph{enumi})}
\tightlist
\item
  \(\mathfrak{q}\) is the normalizer of \(\mathfrak{n}\) in \(\mathfrak{g}\) : let \(\mathfrak{p}\) be this normalizer. Let \(x \in \mathfrak{q}\) and \(y \in \mathfrak{n}\) ; for all \(z \in \mathfrak{q}\) , we have \([z, x] \in \mathfrak{q}\) , so
\end{enumerate}

\[
\Phi ([ x, y ], z) = \Phi (y, [ z, x ]) = 0;
\]

in other words, \([x,y] \in n\) . Hence \(q \subset p\) . Since n is an ideal of p consisting of nilpotent endomorphisms of V, P is orthogonal to n with respect to \(\varPhi\) (Chap. I, no. 3, Prop. 4 d)). Since \(\varPhi\) is non-degenerate \(^{4}\) , \(p \subset q\) , hence our assertion.

\begin{enumerate}
\def\labelenumi{\alph{enumi})}
\setcounter{enumi}{1}
\tightlist
\item
  There exists a reductive Lie subalgebra m in \(\mathfrak{gl}(V)\) such that q is the semi-direct product of m and n: let \(\mathfrak{n}_{\mathrm{V}}(\mathfrak{q})\) be the largest ideal of q consisting of nilpotent endomorphisms of V. Then \(\mathfrak{n}_{\mathrm{V}}(\mathfrak{q})\) contains n, and it is orthogonal to q (loc. cit.); hence \(\mathfrak{n} = \mathfrak{n}_{\mathrm{V}}(\mathfrak{q})\) . Moreover, g is reductive in \(\mathfrak{gl}(V)\) by hypothesis, hence decomposable (Chap. VII, §5, no. 1, Prop. 2); since q is the intersection of g with the normalizer of n in \(\mathfrak{gl}(V)\) , it is a decomposable Lie algebra (loc. cit., Cor. 1 of Prop. 3). Thus, our assertion follows from Prop. 7 of Chap. VII, §5, no. 3.
\end{enumerate}

Choose a Cartan subalgebra \(\mathfrak{h}\) of \(\mathfrak{m}\) ; denote by \(\mathfrak{g}_1\) the commutant of \(\mathfrak{h}\) in \(\mathfrak{g}\) , and put \(\mathfrak{q}_1 = \mathfrak{q} \cap \mathfrak{g}_1, \mathfrak{n}_1 = \mathfrak{n} \cap \mathfrak{g}_1\) .

\begin{enumerate}
\def\labelenumi{\alph{enumi})}
\setcounter{enumi}{2}
\tightlist
\item
  The Lie algebras \(\mathfrak{g}_1, \mathfrak{q}_1\) and \(\mathfrak{n}_1\) satisfy the same hypotheses as \(\mathfrak{g}, \mathfrak{q}\) and \(\mathfrak{n}\) : since \(\mathfrak{m}\) is reductive in \(\mathfrak{gl}(\mathrm{V})\) , \(\mathfrak{h}\) is commutative and is composed of semi-simple endomorphisms of V (Chap. VII, §2, no. 4, Cor. 3 of Th. 2). Thus \(\mathfrak{g}_1 = \mathfrak{g}^0(\mathfrak{h})\) is reductive in \(\mathfrak{g}\) (Chap. VII, §1, no. 3, Prop. 11), hence also in \(\mathfrak{gl}(\mathrm{V})\) (Chap. I, §6, no. 6, Cor. 2 of Prop. 7). It is clear that \(\mathfrak{n}_1\) is composed of nilpotent endomorphisms of V. Since \(\mathfrak{h}\) is a subalgebra of \(\mathfrak{q}\) , reductive in \(\mathfrak{gl}(V)\) , the adjoint representation of \(\mathfrak{h}\) on \(\mathfrak{q}\) is semi-simple; by construction, \(\mathfrak{q}_1\) is the set of invariants of \(\mathrm{ad}_{\mathfrak{q}}(\mathfrak{h})\) , so \(\mathfrak{q} = \mathfrak{q}_1 + [\mathfrak{h}, \mathfrak{q}]\) (Chap. I, §3, no. 5, Prop. 6). Since
\end{enumerate}

\[
\Phi (\mathfrak {g} _ {1}, [ \mathfrak {h}, \mathfrak {q} ]) = \Phi ([ \mathfrak {h}, \mathfrak {g} _ {1} ], \mathfrak {q}) = 0,
\]

an element of \(\mathfrak{g}_1\) is orthogonal to \(\mathfrak{q}_1\) if and only if it is orthogonal to \(\mathfrak{q}\) ; consequently, \(\mathfrak{n}_1 = \mathfrak{g}_1 \cap \mathfrak{n}\) is the orthogonal complement of \(\mathfrak{q}_1\) in \(\mathfrak{g}_1\) .

\begin{enumerate}
\def\labelenumi{\alph{enumi})}
\setcounter{enumi}{3}
\item
  The Cartan subalgebra \(\mathfrak{h}\) of \(\mathfrak{m}\) is a Cartan subalgebra of \(\mathfrak{g}\) : We have \(\mathfrak{q} = \mathfrak{m} \oplus \mathfrak{n}\) and \(\mathfrak{h} = \mathfrak{m} \cap \mathfrak{g}_1\) , so it is immediate that \(\mathfrak{q}_1 = \mathfrak{h} \oplus \mathfrak{n}_1\) . Moreover, \([\mathfrak{h}, \mathfrak{n}_1] = 0\) , \(\mathfrak{h}\) is commutative and \(\mathfrak{n}_1\) is nilpotent, so the Lie algebra \(\mathfrak{q}_1\) is nilpotent. By \(a)\) and \(c)\) , \(\mathfrak{q}_1\) is the normalizer of \(\mathfrak{n}_1\) in \(\mathfrak{g}_1\) ; a fortiori, \(\mathfrak{q}_1\) is equal to its normalizer in \(\mathfrak{g}_1\) , hence is a Cartan subalgebra of \(\mathfrak{g}_1\) . Since \(\mathfrak{g}_1\) is reductive in \(\mathfrak{gl}(V)\) , it follows from Cor. 3 of Th. 2 of Chap. VII, §2, no. 4, that \(\mathfrak{q}_1\) is composed of semi-simple endomorphisms of V; thus, since \(\mathfrak{n}_1\) is composed of nilpotent endomorphisms of V, we have \(\mathfrak{n}_1 = 0\) . Consequently, \(\mathfrak{h} = \mathfrak{q}_1\) is a Cartan subalgebra of \(\mathfrak{g}_1\) , and since \(\mathfrak{g}_1\) normalizes \(\mathfrak{h}\) , we have \(\mathfrak{h} = \mathfrak{g}_1\) . Thus, we have proved that every element of \(\mathfrak{h}\) is a semi-simple element of \(\mathfrak{g}\) , and that the commutant of \(\mathfrak{h}\) in \(\mathfrak{g}\) is equal to \(\mathfrak{h}\) ; it follows that \(\mathfrak{h} = \mathfrak{g}^0(\mathfrak{h})\) , so \(\mathfrak{h}\) is a Cartan subalgebra of \(\mathfrak{g}\) .
\item
  \(\mathfrak{q}\) is a parabolic subalgebra of \(\mathfrak{g}\) : by the preceding, \(\mathfrak{h}\) is a Cartan subalgebra of \(\mathfrak{g}\) , \(\mathfrak{n}\) consists of nilpotent elements of \(\mathfrak{g}\) , and \([\mathfrak{h},\mathfrak{n}]\subset \mathfrak{n}\) . Let \(\bar{k}\) be an algebraic closure of \(k\) ; by definition, \(\mathfrak{q}\) is parabolic in \(\mathfrak{g}\) if and only if \(\bar{k}\otimes_k\mathfrak{q}\) is a parabolic subalgebra of \(\bar{k}\otimes_k\mathfrak{g}\) . The properties stated above being preserved by extension of scalars, for the proof we can restrict ourselves to the case in which \(\mathfrak{h}\) is splitting. Let R be the root system of \((\mathfrak{g},\mathfrak{h})\) ; by Prop. 2 (v) of §3, no. 1, there exists a subset P of R such that \(P\cap (-P) = \varnothing\) and \(\mathfrak{n} = \sum_{\alpha \in P}\mathfrak{g}^{\alpha}\) . Let \(P'\) be the set of roots \(\alpha\) such that \(-\alpha\notin P\) ; we have \(P' \cup (-P') = R\) , and the orthogonal complement \(\mathfrak{q}\) of \(\mathfrak{n}\) in \(\mathfrak{g}\) is equal to \(\mathfrak{h} + \sum_{\alpha \in P'}\mathfrak{g}^{\alpha}\) . We have proved that \(\mathfrak{q}\) is parabolic. Q.E.D.
\end{enumerate}

Lemma 1. Let g be a semi-simple Lie algebra, V a finite dimensional vector space, \(\rho\) a linear representation of g on V, D a vector subspace of V, h a Cartan subalgebra of g, s (resp. \(s'\) ) the set of \(x \in h\) such that \(\rho(x)\mathrm{D} \subset \mathrm{D}\) (resp. \(\rho(x)\mathrm{D} = 0\) ), and \(\Phi\) the bilinear form on g associated \(^{5}\) to \(\rho\) .

\begin{enumerate}
\def\labelenumi{(\roman{enumi})}
\item
  If \(\mathfrak{h}\) is splitting, the vector subspaces \(\mathfrak{s}\) and \(\mathfrak{s}'\) of \(\mathfrak{h}\) are rational over \(\mathbf{Q}\) .
\item
  If \(\rho\) is injective, the restriction of \(\Phi\) to \(\mathfrak{s}\) (resp. \(\mathfrak{s}'\) ) is non-degenerate.
\end{enumerate}

Assume that the Cartan subalgebra \(\mathfrak{h}\) is splitting. Let \(d\) be the dimension of D; put \(\mathrm{W} = \bigwedge^{d}(\mathrm{V})\) and \(\sigma = \bigwedge^{d}(\rho)\) ; denote also by \((e_1, \ldots, e_d)\) a basis of D and \(e = e_1 \wedge \cdots \wedge e_d\) a decomposable \(d\) -vector associated to D. Let P be the set of weights of \(\sigma\) with respect to \(\mathfrak{h}\) ; denote by \(\mathrm{W}^{\mu}\) the subspace of

W associated to the weight \(\mu\) , and put \(e = \sum_{\mu \in P} e^{\mu}\) (with \(e^{\mu} \in W^{\mu}\) for all \(\mu \in P\) ); finally, let \(P'\) be the set of weights \(\mu\) such that \(e^{\mu} \neq 0\) and let \(P''\) be the set of differences of elements of \(P'\) . Let x be in h; then x belongs to s if and only if there exists c in k such that \(\rho(x).e = c.e\) (Chap. VII, §5, no. 4, Lemma 2 (i)). Since \(\rho(x).e^{\mu} = \mu(x).e^{\mu}\) , we see that \(x \in s\) is equivalent to the relation `` \(\mu(x) = 0\) for all \(\mu \in P''\) ''. Now, the Q-structure of h is the Q-vector subspace \(h_{Q}\) of h generated by the coroots \(H_{\alpha}\) and all \(\mu\) in \(P''\) take rational values on \(h_{Q}\) ; it follows (Algebra, Chap. II, §8, no. 4, Prop. 5) that s is a subspace of h rational over Q.

For any weight \(\mu\in P\) , let \(p_{\mu}\) be the projection onto \(V^{\mu}\) associated to the decomposition \(V=\bigoplus_{\mu\in P}V^{\mu}\) ; denote by \(P_{1}\) the set of \(\mu\in P\) such that \(p_{\mu}(D)\neq0\) . It is immediate that \(s'\) is the intersection of the kernels (in h) of the elements of \(P_{1}\) ; it follows, in the same way as for s, that \(s'\) is a subspace of h rational over Q. This proves (i).

By extension of scalars, it suffices to prove (ii) when \(k\) is algebraically closed, hence when \(\mathfrak{h}\) is splitting. Let \(\mathfrak{m}\) be a vector subspace of \(\mathfrak{h}\) rational over \(\mathbf{Q}\) ; for all non-zero \(x\) in \(\mathfrak{m}_{\mathbf{Q}} = \mathfrak{m} \cap \mathfrak{h}_{\mathbf{Q}}\) , we have \(\varPhi(x, x) > 0\) by the Cor. of Prop. 1 of §7, no. 1. The restriction of \(\varPhi\) to \(\mathfrak{m}_{\mathbf{Q}}\) is non-degenerate, and hence so is the restriction of \(\varPhi\) to \(\mathfrak{m}\) since \(\mathfrak{m}\) is canonically isomorphic to \(k \otimes_{\mathbf{Q}} \mathfrak{m}_{\mathbf{Q}}\) .

DEFINITION 1. Let q be a Lie subalgebra of the semi-simple Lie algebra g. Then q is said to be decomposable in g if, for all \(x \in q\) , the semi-simple and nilpotent components of x in g belong to q. Denote by \(\mathfrak{n}_{\mathfrak{g}}(\mathfrak{q})\) the set of elements x of the radical of q such that \(ad_{g}x\) is nilpotent.

Let \(\rho\) be an injective representation of \(\mathfrak{g}\) on a finite dimensional vector space V. We know (Chap. I, §6, no. 3, Th. 3) that an element \(x\) of \(\mathfrak{g}\) is semi-simple (resp. nilpotent) if and only if the endomorphism \(\rho(x)\) of V is semi-simple (resp. nilpotent). It follows immediately that the algebra \(\mathfrak{q}\) is decomposable in \(\mathfrak{g}\) if and only if \(\rho(\mathfrak{q})\) is a decomposable subalgebra of \(\mathfrak{gl}(V)\) in the sense of Definition 1 of Chap. VII, §5, no. 1. With the notations of Chap. VII, §5, no. 3, we also have

\[
\rho (\mathfrak {n} _ {\mathfrak {g}} (\mathfrak {q})) = \mathfrak {n} _ {\mathrm{V}} (\rho (\mathfrak {q})).
\]

THEOREM 2. Let g be a semi-simple Lie algebra, n a subalgebra of g consisting of nilpotent elements, q the normalizer of n in g. Assume that n is the set of nilpotent elements of the radical of q. Then q is parabolic.

Note first of all that q is decomposable (Chap. VII, §5, no. 1, Cor. 1 of Prop. 3). By Th. 1, it suffices to prove that q is the orthogonal complement \(n^{0}\) of n with respect to the Killing form \(\varPhi\) of g. We know that \(q \subset n^{0}\) (Chap. I, §4, no. 3, Prop. 4 d)). By Chap. VII, §5, no. 3, Prop. 7, there exists a subalgebra m of q, reductive in g, such that q is the semi-direct product of m and n.~We show that the restriction of \(\varPhi\) to m is non-degenerate. Let c be the centre of m. We have \(\Phi([\mathfrak{m},\mathfrak{m}],\mathfrak{c})=0\) by Chap. I, §5, no. 5, Prop. 5, and the restriction of \(\Phi\) to \([m,m]\) is non-degenerate by Chap. I, §6, no. 1, Prop. 1. It remains to see that the restriction of \(\Phi\) to c is non-degenerate. Let k be a Cartan subalgebra of \([m,m]\) ; then \(k\oplus c\) is commutative and reductive in g. Let h be a Cartan subalgebra of g containing \(k\oplus c\) (Chap. VII, §2, no. 3, Prop. 10). Then \(h\cap q\) is a commutative subalgebra of q containing \(k\oplus c\) , and \(ad_{q}x\) is semi-simple for all \(x\in h\cap q\) ; hence \(h\cap q\) is contained in a Cartan subalgebra \(h'\) of q (Chap. VII, §2, no. 3, Prop. 10); let f be the projection of q onto m with kernel n; then \(f(h')\) is a Cartan subalgebra of m (Chap. VII, §2, no. 1, Cor. 2 of Prop. 4) containing \(k\oplus c\) , and consequently equal to \(k\oplus c\) ; this proves that \(f(h\cap q)=k\oplus c\) , and since every element of h is semi-simple in g, we have \(h\cap q=k\oplus c\) . Thus,

\[
\mathfrak {c} = \{x \in \mathfrak {h} \mid [ x, \mathfrak {n} ] \subset \mathfrak {n} \text {   and   } [ x, [ \mathfrak {m}, \mathfrak {m} ] ] = 0 \}.
\]

By Lemma 1, the restriction of \(\varPhi\) to \(\mathfrak{c}\) is non-degenerate.

Let \(\mathfrak{q}^0\) be the orthogonal complement of \(\mathfrak{q}\) in \(\mathfrak{g}\) relative to \(\varPhi\) . The preceding proves that \(\mathfrak{q} \cap \mathfrak{q}^0 = \mathfrak{n}\) . Assume that \(\mathfrak{q} \neq \mathfrak{q}^0\) , so \(\mathfrak{q}^0 \neq \mathfrak{n}\) (and \(\mathfrak{q}^0 \supset \mathfrak{n}\) ). Since \(\operatorname{ad}_{\mathfrak{g}} \mathfrak{n}\) leaves \(\mathfrak{q}\) stable, \(\operatorname{ad}_{\mathfrak{g}} \mathfrak{n}\) leaves \(\mathfrak{q}^0\) stable; Engel's theorem proves that there exists \(x \in \mathfrak{q}^0\) such that \(x \notin \mathfrak{n}\) and \([x, \mathfrak{n}] \subset \mathfrak{n}\) . But then \(x \in \mathfrak{q}^0 \cap \mathfrak{q} = \mathfrak{n}\) , a contradiction. Hence \(\mathfrak{q} = \mathfrak{n}^0\) .

COROLLARY 1. Let \(\mathfrak{q}\) be a maximal element of the set of subalgebras of \(\mathfrak{g}\) distinct from \(\mathfrak{g}\) . Then \(\mathfrak{q}\) is either parabolic or reductive in \(\mathfrak{g}\) .

We can assume that g is a Lie subalgebra of \(\mathfrak{gl}(V)\) for some finite dimensional vector space V. Let \(\mathfrak{e}(\mathfrak{q}) \subset \mathfrak{g}\) be the decomposable envelope of q. If \(\mathfrak{e}(\mathfrak{q}) = \mathfrak{g}\) , q is an ideal of g (Chap. VII, §5, no. 2, Prop. 4), hence is semi-simple, and consequently q is reductive in g. Assume that \(\mathfrak{e}(\mathfrak{q}) \neq \mathfrak{g}\) . Then \(\mathfrak{e}(\mathfrak{q}) = \mathfrak{q}\) , so q is decomposable. Assume that q is not reductive in g. Let n be the set of nilpotent elements of the radical of q. Then \(n \neq 0\) (Chap. VII, §5, no. 3, Prop. 7 (i)). Let p be the normalizer of n in g. Then \(p \supset q\) , and \(p \neq g\) since g is semi-simple. Hence p = q. Thus q is parabolic (Th. 1).

COROLLARY 2. Let \(\mathfrak{n}\) be a subalgebra of \(\mathfrak{g}\) consisting of nilpotent elements. There exists a parabolic subalgebra \(\mathfrak{q}\) of \(\mathfrak{g}\) with the following properties:

\begin{enumerate}
\def\labelenumi{(\roman{enumi})}
\item
  \(\mathfrak{n}\subset \mathfrak{n}_{\mathfrak{g}}(\mathfrak{q});\)
\item
  the normalizer of \(\mathfrak{n}\) in \(\mathfrak{g}\) is contained in \(\mathfrak{q}\) ;
\item
  every automorphism of g leaving n invariant leaves q invariant.
\end{enumerate}

If \(\mathfrak{g}\) is splittable, \(\mathfrak{n}\) is contained in a Borel subalgebra of \(\mathfrak{g}\) .

Let \(\mathfrak{q}_1\) be the normalizer of \(\mathfrak{n}\) in \(\mathfrak{g}\) . This is a decomposable subalgebra of \(\mathfrak{g}\) . Let \(\mathfrak{n}_1 = \mathfrak{n}_{\mathfrak{g}}(\mathfrak{q}_1)\) . Define inductively \(\mathfrak{q}_i\) to be the normalizer of \(\mathfrak{n}_{i-1}\) in \(\mathfrak{g}\) , and \(\mathfrak{n}_i\) to be equal to \(\mathfrak{n}_{\mathfrak{g}}(\mathfrak{q}_i)\) . The sequences \((\mathfrak{n}, \mathfrak{n}_1, \mathfrak{n}_2, \ldots)\) and \((\mathfrak{q}_1, \mathfrak{q}_2, \ldots)\) are increasing. There exists \(j\) such that \(\mathfrak{q}_j = \mathfrak{q}_{j+1}\) , in other words \(\mathfrak{q}_j\) is the normalizer of \(\mathfrak{n}_{\mathfrak{g}}(\mathfrak{q}_j)\) in \(\mathfrak{g}\) . Thus \(\mathfrak{q}_j\) is parabolic (Th. 1). We have \(\mathfrak{n} \subset \mathfrak{n}_j =\) \(\mathfrak{n}_{\mathfrak{g}}(\mathfrak{q}_{j})\) , and \(q_{1} \subset q_{j}\) ; every automorphism of g leaving n invariant evidently leaves \(n_{1}, n_{2}, \ldots\) and \(q_{1}, q_{2}, \ldots\) invariant. If g is splittable, \(q_{j}\) contains a Borel subalgebra b, and consequently (§3, no. 4, Prop. 13), we have \(\mathfrak{b} \supset \mathfrak{n}_{\mathfrak{g}}(\mathfrak{q}_{j}) \supset \mathfrak{n}\) .

THEOREM 3. Assume that k is algebraically closed. Let g be a semi-simple Lie algebra. Let a be a solvable subalgebra of g. There exists a Borel subalgebra of g containing a.

By Chap. VII, §5, no. 2, Cor. 1 (ii) of Prop. 4, we can assume that \(\mathfrak{a}\) is decomposable. There exists a commutative subalgebra \(\mathfrak{t}\) of \(\mathfrak{g}\) , consisting of semi-simple elements, such that \(\mathfrak{a}\) is the semi-direct product of \(\mathfrak{t}\) and \(\mathfrak{n}_{\mathfrak{g}}(\mathfrak{a})\) (Chap. VII, §5, no. 3, Cor. 2 of Prop. 6). There exists (Cor. 2 of Th. 2) a parabolic subalgebra \(\mathfrak{q}\) of \(\mathfrak{g}\) such that \(\mathfrak{n}_{\mathfrak{g}}(\mathfrak{a}) \subset \mathfrak{n}_{\mathfrak{g}}(\mathfrak{q})\) , and such that the normalizer of \(\mathfrak{n}_{\mathfrak{g}}(\mathfrak{a})\) in \(\mathfrak{g}\) is contained in \(\mathfrak{q}\) ; a fortiori, \(\mathfrak{a} \subset \mathfrak{q}\) . Let \(\mathfrak{b}\) be a Borel subalgebra of \(\mathfrak{g}\) contained in \(\mathfrak{q}\) and \(\mathfrak{h}\) a Cartan subalgebra of \(\mathfrak{g}\) contained in \(\mathfrak{b}\) . Then \(\mathfrak{h}\) is a Cartan subalgebra of \(\mathfrak{q}\) , so there exists \(s \in \mathrm{Aut}_e(\mathfrak{q})\) such that \(s(\mathfrak{t}) \subset \mathfrak{h}\) (Chap. VII, §2, no. 3, Prop. 10 and Chap. VII, §3, no. 2, Th. 1). We have \(s(\mathfrak{n}_{\mathfrak{g}}(\mathfrak{q})) = \mathfrak{n}_{\mathfrak{g}}(\mathfrak{q})\) (Chap. VII, §3, no. 1, Remark 1), so

\[
s (\mathfrak {a}) = s (t) + s \left(\mathfrak {n} _ {\mathfrak {g}} (\mathfrak {a})\right) \subset \mathfrak {h} + s \left(\mathfrak {n} _ {\mathfrak {g}} (\mathfrak {q})\right) = \mathfrak {h} + \mathfrak {n} _ {\mathfrak {g}} (\mathfrak {q}) \subset \mathfrak {b}.
\]

COROLLARY. If \(k\) is algebraically closed, every maximal solvable subalgebra of \(\mathfrak{g}\) is a Borel subalgebra.

